\subsection{Reproducibility Details}
\label{sec:supplementary-reproducibility}

The CPU-side evaluation was run with Python 3.14.6 on Windows 11. The principal package versions were NumPy 2.5.1, pandas 3.0.5, scikit-learn 1.9.0, SciPy 1.18.0, Biopython 1.87, PyYAML 6.0.3, Matplotlib 3.11.1, seaborn 0.13.2, Jinja2 3.1.6, and pytest 9.1.1. Frozen ESM-2 embeddings were generated in a separate CUDA environment with Python 3.12.13, PyTorch 2.11.0+cu128, model \texttt{facebook/esm2\_t6\_8M\_UR50D}, CUDA, batch size 16, maximum sequence length 1022, 320-dimensional float32 outputs, and 9855 sequences.

The analysis used \texttt{config/default.yaml}. The raw input file \texttt{data/raw/pHopt\_data.csv} was obtained from Zenodo record \texttt{10.5281/zenodo.14252615}, version 1, under CC BY 4.0, and accessed on 2026-10-22 \citep{gado2023phopt}. Its published MD5 is \texttt{d5d0887710e70f2aa5153c15178a6d08}; its local SHA-256 checksum is \texttt{6C2CE7F68E28E4130904A96FC071DBF7FE78E090BD1E408408038F01FA75A3E9}; the cleaned dataset has checksum \texttt{7413D3F48E8A097C4687BA0621F9E4E1D237682FCCAA54B41B952C16E606B3A7}.

Source-label auditing followed the data-quality concerns raised by Zhou and Fang \cite{zhou2026caution}. The named accession B5LXD7 occurs once in the raw file with pHopt 11.5, organism \textit{Bispora sp.}, and EC 3.2.1.78; it occurs in the saved high-error predictions and is a test record for the 50\% cluster-safe and EC-held-out splits. The audit did not overwrite the published label. Instead, a leave-one-flagged-record-out sensitivity analysis was generated by \texttt{scripts/audit\_source\_labels.py}; the accession audit is in \texttt{reports/source\_label\_audit.csv} and the sensitivity results are in \texttt{outputs/metrics/source\_label\_sensitivity.csv}. The original Bispora study is cited as Luo et al. \cite{luo2009novel}.

\subsection{Additional Tables and Figures}
\label{sec:supplementary-additional-results}

Generated supplementary outputs include the complete model-comparison table, the pH-range performance table, bootstrap confidence intervals, paired feature-group comparisons with Holm--Bonferroni-adjusted Wilcoxon p-values, subgroup summaries, and focused error tables under \texttt{outputs/}.

The exhaustive model-comparison table is reproduced below; the main text uses a compact table with readable labels and uncertainty intervals.
\begin{longtable}{lllrrrrrr}
\caption{Test-set performance by feature group, split strategy, and model.} \label{tab:model-comparison} \\
\toprule
feature\_group & split\_strategy & model & n\_examples & mae & rmse & r2 & spearman & pearson \\
\midrule
\endfirsthead
\caption[]{Test-set performance by feature group, split strategy, and model.} \\
\toprule
feature\_group & split\_strategy & model & n\_examples & mae & rmse & r2 & spearman & pearson \\
\midrule
\endhead
\midrule
\multicolumn{9}{r}{Continued on next page} \\
\midrule
\endfoot
\bottomrule
\endlastfoot
esm2\_embedding & cluster\_safe\_mmseqs\_identity\_50 & hist\_gradient\_boosting & 1544 & 0.682200 & 0.998100 & 0.310900 & 0.479300 & 0.561700 \\
esm2\_embedding & cluster\_safe\_mmseqs\_identity\_50 & random\_forest & 1544 & 0.685700 & 1.006900 & 0.298700 & 0.475000 & 0.554100 \\
esm2\_embedding & cluster\_safe\_mmseqs\_identity\_50 & ridge\_regression & 1544 & 0.725100 & 1.022400 & 0.277000 & 0.434100 & 0.529900 \\
handcrafted & cluster\_safe\_mmseqs\_identity\_50 & random\_forest & 1544 & 0.737600 & 1.063900 & 0.217000 & 0.381900 & 0.472000 \\
handcrafted & cluster\_safe\_mmseqs\_identity\_50 & hist\_gradient\_boosting & 1544 & 0.741100 & 1.067100 & 0.212300 & 0.395500 & 0.467200 \\
handcrafted & cluster\_safe\_mmseqs\_identity\_50 & ridge\_regression & 1544 & 0.773900 & 1.106900 & 0.152400 & 0.371000 & 0.398300 \\
esm2\_embedding & cluster\_safe\_mmseqs\_identity\_50 & median\_predictor & 1544 & 0.823800 & 1.253700 & -0.087300 & NaN & NaN \\
handcrafted & cluster\_safe\_mmseqs\_identity\_50 & median\_predictor & 1544 & 0.823800 & 1.253700 & -0.087300 & NaN & NaN \\
esm2\_embedding & cluster\_safe\_mmseqs\_identity\_50 & mean\_predictor & 1544 & 0.851000 & 1.208600 & -0.010400 & NaN & NaN \\
handcrafted & cluster\_safe\_mmseqs\_identity\_50 & mean\_predictor & 1544 & 0.851000 & 1.208600 & -0.010400 & NaN & NaN \\
esm2\_embedding & cluster\_safe\_mmseqs\_identity\_70 & hist\_gradient\_boosting & 1481 & 0.628200 & 0.883700 & 0.407000 & 0.461300 & 0.638700 \\
esm2\_embedding & cluster\_safe\_mmseqs\_identity\_70 & random\_forest & 1481 & 0.629800 & 0.893700 & 0.393500 & 0.473800 & 0.631800 \\
esm2\_embedding & cluster\_safe\_mmseqs\_identity\_70 & ridge\_regression & 1481 & 0.682700 & 0.923100 & 0.353000 & 0.422800 & 0.596400 \\
handcrafted & cluster\_safe\_mmseqs\_identity\_70 & random\_forest & 1481 & 0.695400 & 0.984800 & 0.263600 & 0.375900 & 0.515400 \\
handcrafted & cluster\_safe\_mmseqs\_identity\_70 & hist\_gradient\_boosting & 1481 & 0.701400 & 0.979100 & 0.272100 & 0.376600 & 0.522300 \\
handcrafted & cluster\_safe\_mmseqs\_identity\_70 & ridge\_regression & 1481 & 0.749900 & 1.042800 & 0.174300 & 0.352000 & 0.420800 \\
esm2\_embedding & cluster\_safe\_mmseqs\_identity\_70 & median\_predictor & 1481 & 0.781800 & 1.172600 & -0.044100 & NaN & NaN \\
handcrafted & cluster\_safe\_mmseqs\_identity\_70 & median\_predictor & 1481 & 0.781800 & 1.172600 & -0.044100 & NaN & NaN \\
esm2\_embedding & cluster\_safe\_mmseqs\_identity\_70 & mean\_predictor & 1481 & 0.812800 & 1.148000 & -0.000600 & NaN & NaN \\
handcrafted & cluster\_safe\_mmseqs\_identity\_70 & mean\_predictor & 1481 & 0.812800 & 1.148000 & -0.000600 & NaN & NaN \\
esm2\_embedding & cluster\_safe\_mmseqs\_identity\_90 & hist\_gradient\_boosting & 1473 & 0.614800 & 0.868300 & 0.414400 & 0.513300 & 0.644500 \\
esm2\_embedding & cluster\_safe\_mmseqs\_identity\_90 & random\_forest & 1473 & 0.616400 & 0.881900 & 0.395900 & 0.515400 & 0.633000 \\
esm2\_embedding & cluster\_safe\_mmseqs\_identity\_90 & ridge\_regression & 1473 & 0.678600 & 0.924800 & 0.335700 & 0.450600 & 0.582900 \\
handcrafted & cluster\_safe\_mmseqs\_identity\_90 & random\_forest & 1473 & 0.681300 & 0.959700 & 0.284600 & 0.402900 & 0.535300 \\
handcrafted & cluster\_safe\_mmseqs\_identity\_90 & hist\_gradient\_boosting & 1473 & 0.705200 & 0.976600 & 0.259200 & 0.359100 & 0.511700 \\
handcrafted & cluster\_safe\_mmseqs\_identity\_90 & ridge\_regression & 1473 & 0.756000 & 1.035100 & 0.167900 & 0.321700 & 0.411200 \\
esm2\_embedding & cluster\_safe\_mmseqs\_identity\_90 & median\_predictor & 1473 & 0.781900 & 1.155600 & -0.037300 & NaN & NaN \\
handcrafted & cluster\_safe\_mmseqs\_identity\_90 & median\_predictor & 1473 & 0.781900 & 1.155600 & -0.037300 & NaN & NaN \\
esm2\_embedding & cluster\_safe\_mmseqs\_identity\_90 & mean\_predictor & 1473 & 0.817800 & 1.136400 & -0.003000 & NaN & NaN \\
handcrafted & cluster\_safe\_mmseqs\_identity\_90 & mean\_predictor & 1473 & 0.817800 & 1.136400 & -0.003000 & NaN & NaN \\
esm2\_embedding & dataset\_provided & random\_forest & 1971 & 0.603500 & 0.849100 & 0.459200 & 0.598800 & 0.691300 \\
esm2\_embedding & dataset\_provided & hist\_gradient\_boosting & 1971 & 0.613200 & 0.857200 & 0.448800 & 0.574800 & 0.672300 \\
handcrafted & dataset\_provided & random\_forest & 1971 & 0.672700 & 0.934300 & 0.345100 & 0.483100 & 0.595100 \\
esm2\_embedding & dataset\_provided & ridge\_regression & 1971 & 0.674200 & 0.912500 & 0.375300 & 0.497300 & 0.613800 \\
handcrafted & dataset\_provided & hist\_gradient\_boosting & 1971 & 0.688300 & 0.944500 & 0.330700 & 0.451100 & 0.578100 \\
handcrafted & dataset\_provided & ridge\_regression & 1971 & 0.757300 & 1.020300 & 0.219100 & 0.380900 & 0.472800 \\
esm2\_embedding & dataset\_provided & median\_predictor & 1971 & 0.809700 & 1.181700 & -0.047600 & NaN & NaN \\
handcrafted & dataset\_provided & median\_predictor & 1971 & 0.809700 & 1.181700 & -0.047600 & NaN & NaN \\
esm2\_embedding & dataset\_provided & mean\_predictor & 1971 & 0.848100 & 1.155200 & -0.001100 & NaN & NaN \\
handcrafted & dataset\_provided & mean\_predictor & 1971 & 0.848100 & 1.155200 & -0.001100 & NaN & NaN \\
esm2\_embedding & ec\_held\_out & random\_forest & 1478 & 0.832100 & 1.188800 & 0.233300 & 0.511200 & 0.496400 \\
esm2\_embedding & ec\_held\_out & ridge\_regression & 1478 & 0.839500 & 1.161400 & 0.268200 & 0.527500 & 0.525100 \\
esm2\_embedding & ec\_held\_out & hist\_gradient\_boosting & 1478 & 0.840900 & 1.195200 & 0.225100 & 0.501700 & 0.494400 \\
handcrafted & ec\_held\_out & random\_forest & 1478 & 0.884400 & 1.211800 & 0.203400 & 0.479200 & 0.483300 \\
handcrafted & ec\_held\_out & hist\_gradient\_boosting & 1478 & 0.913900 & 1.241200 & 0.164300 & 0.437600 & 0.456100 \\
handcrafted & ec\_held\_out & ridge\_regression & 1478 & 0.933900 & 1.257000 & 0.142900 & 0.470400 & 0.454800 \\
esm2\_embedding & ec\_held\_out & mean\_predictor & 1478 & 1.099700 & 1.486500 & -0.198800 & NaN & NaN \\
handcrafted & ec\_held\_out & mean\_predictor & 1478 & 1.099700 & 1.486500 & -0.198800 & NaN & NaN \\
esm2\_embedding & ec\_held\_out & median\_predictor & 1478 & 1.134300 & 1.551900 & -0.306600 & NaN & NaN \\
handcrafted & ec\_held\_out & median\_predictor & 1478 & 1.134300 & 1.551900 & -0.306600 & NaN & NaN \\
esm2\_embedding & exact\_duplicate\_safe & hist\_gradient\_boosting & 1480 & 0.602800 & 0.862500 & 0.440300 & 0.560000 & 0.665000 \\
esm2\_embedding & exact\_duplicate\_safe & random\_forest & 1480 & 0.607300 & 0.871700 & 0.428200 & 0.562600 & 0.662200 \\
handcrafted & exact\_duplicate\_safe & random\_forest & 1480 & 0.663800 & 0.944300 & 0.329000 & 0.482000 & 0.578400 \\
handcrafted & exact\_duplicate\_safe & hist\_gradient\_boosting & 1480 & 0.679700 & 0.953700 & 0.315600 & 0.463400 & 0.562900 \\
esm2\_embedding & exact\_duplicate\_safe & ridge\_regression & 1480 & 0.686300 & 0.943600 & 0.330000 & 0.470000 & 0.577600 \\
handcrafted & exact\_duplicate\_safe & ridge\_regression & 1480 & 0.751800 & 1.036900 & 0.190900 & 0.379400 & 0.439900 \\
esm2\_embedding & exact\_duplicate\_safe & median\_predictor & 1480 & 0.798400 & 1.179100 & -0.046100 & NaN & NaN \\
handcrafted & exact\_duplicate\_safe & median\_predictor & 1480 & 0.798400 & 1.179100 & -0.046100 & NaN & NaN \\
esm2\_embedding & exact\_duplicate\_safe & mean\_predictor & 1480 & 0.837100 & 1.153200 & -0.000700 & NaN & NaN \\
handcrafted & exact\_duplicate\_safe & mean\_predictor & 1480 & 0.837100 & 1.153200 & -0.000700 & NaN & NaN \\
esm2\_embedding & organism\_held\_out & random\_forest & 1478 & 0.536000 & 0.758500 & 0.359000 & 0.457700 & 0.605500 \\
esm2\_embedding & organism\_held\_out & hist\_gradient\_boosting & 1478 & 0.538600 & 0.754100 & 0.366400 & 0.455300 & 0.606200 \\
handcrafted & organism\_held\_out & random\_forest & 1478 & 0.594400 & 0.837400 & 0.218700 & 0.367100 & 0.474100 \\
handcrafted & organism\_held\_out & hist\_gradient\_boosting & 1478 & 0.605200 & 0.842600 & 0.208900 & 0.335500 & 0.464500 \\
esm2\_embedding & organism\_held\_out & ridge\_regression & 1478 & 0.620900 & 0.835200 & 0.222700 & 0.386800 & 0.502600 \\
esm2\_embedding & organism\_held\_out & median\_predictor & 1478 & 0.630600 & 0.947400 & -0.000100 & NaN & NaN \\
handcrafted & organism\_held\_out & median\_predictor & 1478 & 0.630600 & 0.947400 & -0.000100 & NaN & NaN \\
handcrafted & organism\_held\_out & ridge\_regression & 1478 & 0.659300 & 0.907600 & 0.082200 & 0.286200 & 0.345800 \\
esm2\_embedding & organism\_held\_out & mean\_predictor & 1478 & 0.709900 & 0.973900 & -0.056900 & NaN & NaN \\
handcrafted & organism\_held\_out & mean\_predictor & 1478 & 0.709900 & 0.973900 & -0.056900 & NaN & NaN \\
esm2\_embedding & random & random\_forest & 1479 & 0.625000 & 0.911400 & 0.407600 & 0.539300 & 0.645900 \\
esm2\_embedding & random & hist\_gradient\_boosting & 1479 & 0.633500 & 0.907900 & 0.412100 & 0.514900 & 0.643300 \\
handcrafted & random & random\_forest & 1479 & 0.698400 & 1.011400 & 0.270500 & 0.432900 & 0.520300 \\
esm2\_embedding & random & ridge\_regression & 1479 & 0.703500 & 0.975000 & 0.322000 & 0.453300 & 0.567900 \\
handcrafted & random & hist\_gradient\_boosting & 1479 & 0.709100 & 1.004900 & 0.279800 & 0.414600 & 0.530300 \\
handcrafted & random & ridge\_regression & 1479 & 0.781700 & 1.082900 & 0.163700 & 0.333200 & 0.405000 \\
esm2\_embedding & random & median\_predictor & 1479 & 0.812700 & 1.218400 & -0.058700 & NaN & NaN \\
handcrafted & random & median\_predictor & 1479 & 0.812700 & 1.218400 & -0.058700 & NaN & NaN \\
esm2\_embedding & random & mean\_predictor & 1479 & 0.844400 & 1.184300 & -0.000200 & NaN & NaN \\
handcrafted & random & mean\_predictor & 1479 & 0.844400 & 1.184300 & -0.000200 & NaN & NaN \\
\end{longtable}

